\begin{explanationbox}
\textbf{\textcolor{CExplanation}{一句话直觉}}

互为倒数的两个实数的和，其取值范围由它们的符号决定：同号时和不小于 $2$，异号时和不大于 $-2$。

\medskip
\textbf{\textcolor{CExplanation}{核心拆解}}
\begin{description}[style=nextline, leftmargin=2.8em, labelsep=0.5em, itemsep=0.45em]
    \item[\textbf{要点一（符号定方向）：}] 不等式的方向（$\geq$ 或 $\leq$）完全由 $a,b$ 的乘积符号决定。同号时，和有一个下界 $2$；异号时，和有一个上界 $-2$。
    \item[\textbf{要点二（平方非负）：}] 证明的核心是利用平方的非负性：$\bigl(\sqrt{\frac{a}{b}} - \sqrt{\frac{b}{a}}\bigr)^2 \geq 0$ 或 $\bigl(\frac{a}{b} + \frac{b}{a}\bigr)^2 \geq 4$，但需要注意符号处理。
\end{description}

\medskip
\textbf{\textcolor{CExplanation}{几何本质（双曲线图像）}}

考虑函数 $y = x + \frac{1}{x}$ 的图像。当 $x>0$ 时，函数在 $x=1$ 处取得最小值 $2$；当 $x<0$ 时，函数在 $x=-1$ 处取得最大值 $-2$。这直观地展示了结论的几何意义。

\medskip
\textbf{\textcolor{CExplanation}{代数意义（完全平方）}}

代数上，结论源于完全平方公式：$\left(\sqrt{\frac{a}{b}} - \sqrt{\frac{b}{a}}\right)^2 \geq 0$ 展开即得。注意当 $ab>0$ 时，平方根有意义；当 $ab<0$ 时，需稍作变形，例如考虑 $\left(\sqrt{-\frac{a}{b}} - \sqrt{-\frac{b}{a}}\right)^2 \geq 0$。

\medskip
\textbf{\textcolor{CExplanation}{考点价值（最值问题）}}
\begin{itemize}[leftmargin=*, itemsep=0.5em, topsep=0.4em]
    \item \textbf{考法一（直接应用）：} 在选择题或填空题中，直接要求求倒数和的最值。
    \item \textbf{考法二（证明工具）：} 作为中间步骤证明其他不等式，或求解含参数的范围问题。
\end{itemize}

\medskip
\textbf{\textcolor{CExplanation}{顿悟点}}

\textbf{符号决定一切：} 使用前必须先判断 $a,b$ 的符号，否则容易弄错不等号方向。当 $ab>0$ 时，和有一个最小值；当 $ab<0$ 时，和有一个最大值。

\medskip
\textbf{\textcolor{CExplanation}{使用场景}}
\begin{itemize}[leftmargin=*, itemsep=0.5em, topsep=0.4em]
    \item \textbf{场景一（求最值）：} 当表达式中出现倒数和形式，且变量符号已知时，可直接应用此结论求最值。
    \item \textbf{场景二（放缩证明）：} 在不等式证明中，可将倒数和放缩为一个常数，简化问题。
\end{itemize}
\end{explanationbox}
